\chapter{Hamilton方程与相空间}

\section{为什么把速度换成动量}
Lagrange方程以$(q,\dot q)$描述状态，通常是二阶方程。
Hamilton方法改用$(q,p)$，得到一阶方程组，并把守恒律、正则变换和统计力学的相空间联系起来。
这里$p_i$是广义动量，只有在特定坐标和相互作用下才等于$m\dot q_i$。

设$L(q,\dot q,t)$足够光滑，且速度Hessian
\begin{equation}
 W_{ij}=\frac{\p^2L}{\p\dot q_i\p\dot q_j}
\end{equation}
可逆。由隐函数定理，$p_i=\p L/\p\dot q_i$可局部解出$\dot q_i(q,p,t)$。
在这样的正则区域内定义Legendre变换
\begin{equation}
 H(q,p,t)=\sum_i p_i\dot q_i-L(q,\dot q,t).
\end{equation}
右边的速度最后必须用$q,p,t$表示，求偏导时把$q,p$视为独立变量。
求全微分并使用Lagrange方程，
\begin{equation}
 \d H=\sum_i\dot q_i\d p_i-\sum_i\dot p_i\d q_i-\pv L t\d t,
\end{equation}
得到Hamilton方程
\begin{equation}
 \boxed{\dot q_i=\pv H{p_i},\qquad \dot p_i=-\pv H{q_i}},\qquad
 \pv Ht=-\pv Lt.
 \label{eq:cm-hamilton}
\end{equation}
这里$\p H/\p t$固定$q,p$，而$\p L/\p t$固定$q,\dot q$，再将速度代回。
上述推导假定无另加的非保守广义力。速度Hessian奇异时，动量之间可能出现约束，
不能直接求逆；需要另外的约束Hamilton方法。

沿运动轨迹，$\dot H=\p H/\p t$，所以自治Hamiltonian守恒。
但$H$是否等于实验室系机械能，仍需依据坐标选择和相互作用判断。
时间相关的正则变换可以使新Hamiltonian与原能量相差一项。

\begin{example}{电磁场中的正则动量}{cm-magnetic-hamilton}
在SI制中，电荷$e$、质量$m$的非相对论粒子具有
\begin{equation}
 L=\frac12m\bm v^2+e\bm A(\bm r,t)\cdot\bm v-e\phi(\bm r,t).
\end{equation}
由分量求导可恢复$ m\dot{\bm v}=e(\bm E+\bm v\times\bm B)$，其中
$\bm E=-\nabla\phi-\p_t\bm A$，$\bm B=\nabla\times\bm A$。
正则动量与Hamiltonian是
\begin{equation}
 \bm p=m\bm v+e\bm A,\qquad
 H=\frac{|\bm p-e\bm A|^2}{2m}+e\phi.
\end{equation}
因此机械动量$\bm\pi=m\bm v$一般不是正则动量。
在规范变换$\bm A\mapsto\bm A+\nabla\chi$、$\phi\mapsto\phi-\p_t\chi$下，
$L$增加$e\d\chi/\d t$，运动方程不变，$\bm p$则相应改变。
\end{example}

\section{Poisson括号组织运动方程}
对相空间函数$f,g$，采用约定
\begin{equation}
 \{f,g\}=\sum_i\left(\pv f{q_i}\pv g{p_i}-\pv f{p_i}\pv g{q_i}\right).
\end{equation}
于是
\begin{equation}
 \dot f=\{f,H\}+\pv ft,\qquad
 \{q_i,p_j\}=\delta_{ij},\quad \{q_i,q_j\}=\{p_i,p_j\}=0.
\end{equation}
不显含时间的$f$对该区域内所有轨道守恒，当且仅当$\{f,H\}=0$在该区域恒成立。
括号满足反对称性、双线性、乘积规则和Jacobi恒等式；由Jacobi恒等式可证，
两个不显含时间的守恒量的Poisson括号也是守恒量。

在三维Cartesian正则坐标中，$\bm\ell=\bm r\times\bm p$满足
\begin{equation}
 \{\ell_i,\ell_j\}=\sum_k\epsilon_{ijk}\ell_k,\qquad
 \{\ell^2,\ell_i\}=0.
\end{equation}
例如$\ell_z=xp_y-yp_x$生成绕$z$轴的无穷小主动旋转：
$\delta x=\epsilon\{x,\ell_z\}=-\epsilon y$，$\delta y=\epsilon x$。
机械角动量在带电粒子问题中不能未经检查就替代上述正则表达式。
同样，由$\pi_i=p_i-eA_i$可得
\begin{equation}
 \{\pi_i,\pi_j\}=e\sum_k\epsilon_{ijk}B_k,
\end{equation}
说明磁场中机械动量分量并不相互Poisson对易。

\section{Liouville定理：流动不压缩正则体积}
令$\bm z=(q_1,\ldots,q_n,p_1,\ldots,p_n)$，Hamilton方程写成
\begin{equation}
 \dot{\bm z}=J\nabla H,\qquad
 J=\begin{pmatrix}0&\mathbf 1\\-\mathbf 1&0\end{pmatrix}.
\end{equation}
对$C^2$的$H$，流场散度为
\begin{equation}
 \nabla_{\bm z}\cdot\dot{\bm z}
 =\sum_i\left(\frac{\p^2H}{\p q_i\p p_i}
 -\frac{\p^2H}{\p p_i\p q_i}\right)=0.
\end{equation}
因此随轨道流动的小体积$\prod_i\d q_i\d p_i$保持不变，即使$H$显含时间仍成立。
若$\rho(q,p,t)$是粒子系综的概率密度，由连续性方程得
\begin{equation}
 \pv\rho t+\{\rho,H\}=0,
\end{equation}
即密度沿轨道不变。即使$q$是曲线坐标，只要$p$是其共轭动量，正则体积仍为$\prod_i\d q_i\d p_i$；
换成一般非正则相空间变量时，需相应变换体积测度。

对于阻尼谐振子，以$p=m\dot x$写$\dot x=p/m$、$\dot p=-kx-2\gamma p$，
散度为$-2\gamma$，相面积收缩。它不是上述封闭Hamilton系统，
不会构成Liouville定理的反例。
相体积守恒也不表示每个形状保持不变，细长拉伸和折叠完全可能发生。

\section{从解析结构到数值积分}
对可分离Hamiltonian$H=p^2/(2m)+V(q)$，取固定步长$h>0$，速度Verlet方法为
\begin{align}
 p_{n+1/2}&=p_n-\frac h2V'(q_n),\\
 q_{n+1}&=q_n+\frac hm p_{n+1/2},\\
 p_{n+1}&=p_{n+1/2}-\frac h2V'(q_{n+1}).
\end{align}
它由势能与动能各自的精确Hamilton流组合而成，具有二阶精度、时间可逆性和辛结构。
它一般不精确守恒原Hamiltonian，但在适当步长和光滑有界运动中常有较好的长期能量行为。
对频率$\omega_0>0$的谐振子，线性稳定要求$0<h\omega_0<2$；边界$h\omega_0=2$一般已有线性增长。
辛性不能挽救过大的时间步长；若直接把$h$改成依赖当前状态的自适应步长，通常也不再保持辛性。
这提供了实用核对：数值轨道看似平滑仍不够，还应看能量误差、相空间形状及步长收敛。

\paragraph{复习与阅读}
首先检查Legendre变换能否求逆，再区分正则动量、机械动量和能量。
Poisson括号给出演化与对称性，Liouville定理给出相流的体积性质。
延伸阅读：Morin公开的Hamilton补充章；Tong讲义第4章。
